problem:  
Let \( G = SU(n) \) with maximal torus \( T = \exp(\mathfrak{t}) \) and root system \( \Phi \subset \mathfrak{t}^* \).  
For each nonempty subset \( \Phi_0 \subset \Phi^+ \), consider the left-invariant, bracket-generating distribution  
\[
\Delta_{\Phi_0}(g) = L_g\Big( \bigoplus_{\alpha \in \Phi_0} (\mathfrak{g}_\alpha \oplus \mathfrak{g}_{-\alpha}) \Big)
\]
endowed with the sub-Riemannian metric \(g_{\Phi_0}\) induced by the negative Killing form.  
Let \(\mathrm{Cut}_{\Phi_0}(e)\) denote the cut locus at the identity for this sub-Riemannian structure.

**Main open question:**  
Determine for which \( \Phi_0 \subset \Phi^+ \) the cut locus \( \mathrm{Cut}_{\Phi_0}(e) \) is **invariant under the Weyl group action** and lies entirely inside the union of Weyl chamber walls:
\[
\mathrm{Cut}_{\Phi_0}(e) \subseteq \bigcup_{w \in W} w(\exp(\mathfrak{t}_\mathrm{wall})),
\]
where \(W\) is the Weyl group and \(\mathfrak{t}_\mathrm{wall}\subset \mathfrak{t}\) is the union of hyperplanes fixed by nontrivial Weyl reflections.  

Equivalently, one asks:  
> Do there exist left-invariant sub-Riemannian structures on \(SU(n)\) defined by root-space subsets for which all loss-of-optimality phenomena of horizontal geodesics occur along the natural singular strata determined by Weyl symmetry, and no other directions?

Further, for such distributions (if any exist), does the normal Hamiltonian flow admit a moment-map interpretation whose components generate commuting flows invariant under the Weyl action?

---

Why is it a "good" problem:  
This version poses a plausible, geometrically and algebraically consistent research problem that bridges Lie theory, sub-Riemannian geometry, and integrable systems. The question asks whether **the global geometric singularities of the sub-Riemannian exponential map––the cut and conjugate loci––are entirely dictated by Weyl symmetry**, i.e., whether the loss of geodesic optimality is governed purely by algebraic reflection structures intrinsic to \(SU(n)\).  

It is “good” because:

* It frames a **sharp and concrete conjecture** about the qualitative geometry of invariant sub-Riemannian metrics, rooted in natural algebraic data (root systems, Weyl group).
* It **connects analytic control theory (cut/conjugate loci in nonholonomic systems)** with **algebraic symmetry and representation theory**, forging a bridge between geometric control and Lie-theoretic orbit stratification.
* It remains **simple to state yet profoundly deep**: even in \(SU(3)\) the geometry of \(\mathrm{Cut}_{\Phi_0}(e)\) is unknown, so any identification or restriction to Weyl-invariant strata would represent a major conceptual breakthrough.  
* It naturally invites further directions—studying the role of abnormal extremals or moment-map-integrable Hamiltonians—anchored to precise, established mathematical structures.

Hence, the question is mathematically valid, potentially far-reaching, and truly embodies the “narrow entrance, profound depth” character of a good research-level problem.